\chapter{Narcolepsy}

\section*{Definition}
Narcolepsy is a chronic neurological disorder characterized by excessive daytime sleepiness and sudden, uncontrollable episodes of falling asleep. It is caused by loss of hypocretin (orexin) neurons in the hypothalamus.

\section*{What Happens in Your Body}
Narcolepsy type 1 results from autoimmune destruction of hypocretin-producing neurons in the lateral hypothalamus. Hypocretin is a neurotransmitter that promotes wakefulness and stabilizes sleep-wake transitions. Without hypocretin, people experience rapid transitions into REM sleep, causing cataplexy (sudden muscle weakness triggered by emotions), sleep paralysis, and hypnagogic hallucinations.

\section*{Causes}
\begin{itemize}
\item Autoimmune destruction of hypocretin neurons (type 1 narcolepsy)
\item Genetic predisposition (HLA-DQB1*0602 allele)
\item Environmental triggers: influenza A (H1N1) infection or vaccination
\item Traumatic brain injury affecting the hypothalamus
\item Brain tumors or strokes involving the hypothalamus
\item Multiple sclerosis affecting the hypothalamus
\item Type 2 narcolepsy (idiopathic, without cataplexy)
\end{itemize}

\section*{When to See a Doctor}
\begin{itemize}
\item Excessive daytime sleepiness with irresistible sleep attacks
\item Cataplexy (sudden loss of muscle tone triggered by laughter, anger, or surprise)
\item Sleep paralysis (inability to move upon falling asleep or waking)
\item Hypnagogic hallucinations (vivid, dream-like experiences at sleep onset)
\item Fragmented nighttime sleep
\item Automatic behaviors (performing tasks without awareness)
\item Weight gain and obesity
\end{itemize}

\section*{Related Symptoms}
\begin{itemize}
\item Excessive daytime sleepiness
\item Fatigue
\item Insomnia
\item Sleep apnea
\item Depression
\item Hallucinations (hypnagogic)
\item Obesity
\end{itemize}

\section*{Self-Care / Home Management}
\begin{itemize}
\item Schedule short (15-20 minute) naps at regular times
\item Maintain a consistent sleep schedule
\item Avoid alcohol and caffeine close to bedtime
\item Exercise daily (but not near bedtime)
\item Inform employers, teachers, and family about the condition
\item Avoid driving or operating heavy machinery if untreated
\item Join a support group
\end{itemize}

\section*{How Your Doctor Will Evaluate You}
\begin{itemize}
\item Polysomnography (overnight sleep study)
\item Multiple Sleep Latency Test (MSLT) to measure sleep propensity
\item Mean sleep latency less than 8 minutes and 2 or more sleep-onset REM periods
\item HLA typing (presence of DQB1*0602)
\item Measurement of CSF hypocretin-1 levels (less than 110 pg/mL in type 1)
\end{itemize}

\section*{Treatment Options}
\begin{itemize}
\item Stimulants: modafinil, armodafinil (first-line for daytime sleepiness)
\item Sodium oxybate (gamma-hydroxybutyrate) for cataplexy and fragmented sleep
\item Antidepressants: venlafaxine, fluoxetine for cataplexy
\item Pitolisant (histamine H3 receptor inverse agonist)
\item Solriamfetol (dopamine/norepinephrine reuptake inhibitor)
\item Scheduled naps and good sleep hygiene
\end{itemize}

\section*{References}
\begin{thebibliography}{99}
\bibitem{ref1} Scammell TE. "Narcolepsy." \textit{New England Journal of Medicine}, 2015;373(27):2654--2662.
\bibitem{ref2} Kornum BR, Knudsen S, Ollila HM, et al. "Narcolepsy." \textit{Nature Reviews Disease Primers}, 2017;3:16100.
\bibitem{ref3} American Academy of Sleep Medicine. \textit{International Classification of Sleep Disorders}, 3rd edition. American Academy of Sleep Medicine, 2014.
\bibitem{ref4} Mignot E. "Genetics of narcolepsy." \textit{Sleep Medicine Clinics}, 2014;9(2):153--168.
\bibitem{ref5} Thorpy MJ, Dauvilliers Y. "Clinical and practical considerations in the pharmacologic management of narcolepsy." \textit{Sleep Medicine}, 2015;16(1):9--18.
\bibitem{ref6} Billiard M, Dauvilliers Y, Dolenc-Groselj L, et al. "Management of narcolepsy in adults." \textit{Sleep Medicine Reviews}, 2020;52:101304.

\end{thebibliography}

\clearpage
